\chapter{Conclusion and Recommendations}
\label{chap:conclusion}

\section{Conclusion}

This thesis developed and evaluated an interpretable regime-aware model for twelve-month loan-level mortgage probability of default. The empirical pipeline joined Freddie Mac origination and monthly performance records to transformed FRED macroeconomic series, constructed a forward-looking 90-DPD outcome for eligible loan-months, and compressed four adverse-oriented macroeconomic factors into a filtered probability from a two-state Gaussian HMM. A regularised logistic model used that probability as a common level term and interacted it with selected borrower, leverage, interest-rate, and payment-history features. Borrower-only and additive macroeconomic logistic models supplied the principal comparisons, and validation-fitted Platt scaling produced the primary test probabilities.

The first research question concerned the macroeconomic regimes identified by the HMM. The terminal fit assigned the entire April 2023--September 2024 test period to hard stress with probabilities above 0.999902. The causal rolling fits produced both hard states across 36 months, but 35 months still had probabilities at or beyond 0.001 and 0.999. The HMM therefore identified persistent and sharply separated fitted periods. It did not provide enough intermediate posterior variation to justify a smooth stress-gradient interpretation or reliable within-fold hard-regime comparisons.

The second question concerned how regime information changed the PD specification. The final equation added $q_t$ to the loan-level logit and allowed eight prespecified or deterministically selected borrower and payment-history effects to vary with it. This structure created an explicit, auditable macro-to-credit pathway without retaining the full correlated macroeconomic vector in the main regime-aware model. Its interpretation is conditional and predictive: it identifies which sensitivities the model permits to vary, not causal effects of economic stress.

The third question concerned comparative performance. In the original terminal split, regime-aware achieved modestly better proper scores and discrimination than borrower-only. That result was accompanied by weaker validation performance, worse monthly count calibration, and a saturated terminal state. The predefined causal rolling analysis therefore became the more authoritative robustness comparison. Regime-aware won the first rolling origin and lost the next two. Across 6,650,858 pooled rolling-test loan-months, borrower-only had the lower Brier score and log loss, the higher ROC-AUC and PR-AUC, and the smaller monthly default-count error. Paired three-month moving-block intervals included zero, showing that the exact difference is imprecise across only 36 months, while the point estimates and metric ordering favour borrower-only.

The final conclusion is deliberately calibration-first. The regime-aware model is an interpretable, period-dependent diagnostic, not a universally more accurate forecasting model. Borrower-only logistic regression is stronger in pooled causal proper scores, discrimination, and monthly count calibration in this study. Both principal models overpredicted the pooled rolling event rate, so even the preferred comparative model requires ongoing calibration monitoring.

\section{Contributions}

The study makes four bounded contributions. First, it defines a reproducible loan-month estimand whose outcome, horizon, risk set, and information timing are explicit. Second, it links a filtered latent macroeconomic state to loan-level PD through a small and interpretable interaction structure. Third, it evaluates probabilities using proper scores, ranking measures, grouped and mean calibration, and monthly expected-default counts rather than relying on discrimination alone. Fourth, it demonstrates empirically that a favourable terminal result can reverse under causal rolling-origin fitting when upstream macroeconomic transformations and HMM parameters are also restricted to training information.

The fourth contribution is central. The stricter evaluation does not merely add another table; it changes the strength of the conclusion. It prevents a full-window latent-state estimate from being treated as if it were generated in real time, and it makes temporal instability visible without introducing new models or selecting a calibrator from test performance. The mixed result is scientifically useful because it defines both the potential and the boundary of the proposed architecture.

\section{Recommendations}

For the frozen empirical comparison, borrower-only logistic regression should be reported as the primary forecasting model. The regime-aware specification should be retained as an interpretable challenger and diagnostic. The additive macroeconomic model should remain a baseline demonstrating the instability of the tested direct specification; its poor rolling performance should not be generalised to all uses of macroeconomic variables.

Any operational implementation should refit transformations, scaling, HMM parameters, logistic coefficients, and calibration maps using only information available at the forecast origin. Monitoring should report calibration-in-the-large, Brier score, log loss, grouped calibration, monthly predicted and observed default counts, ROC-AUC, and PR-AUC. Because the HMM posterior was saturated, monitoring should also include the share of extreme probabilities, hard-state composition, entropy, transition matrices, and the selected deterministic initialisation.

Future research should prioritise external and temporal validation rather than further optimisation on the frozen tests. A first extension is evaluation on later untouched months. A second is replication on a different mortgage or consumer-credit portfolio, especially an African dataset if the work is positioned for an African journal audience. Such a replication should redefine the outcome and macroeconomic inputs for the local product and institutional setting rather than transfer United States coefficients.

A third extension is a genuine real-time macroeconomic exercise using historical data vintages where available. A fourth is uncertainty analysis that resamples both calendar blocks and borrower clusters or uses a hierarchical design suited to repeated loan-month observations. Alternative state counts, covariance structures, duration models, or Bayesian state uncertainty may be studied prospectively with a new validation protocol. They should not be selected retrospectively from the test periods reported here.

Finally, threshold decisions, economic value, fairness, and governance require separate analysis. The present thesis evaluates probability forecasts. Approval, pricing, collections, and capital applications would need explicit costs, capacity constraints, protected-group assessment, model-risk controls, and locally relevant validation.

\section{Closing Statement}

Macroeconomic regimes offer a coherent way to organise changing credit conditions, but coherence does not guarantee stable forecast gains. In this study the regime-aware structure was transparent and occasionally advantageous, yet its benefit depended on the evaluation period and its posterior signal was nearly binary. The borrower-only model provided the stronger pooled causal forecast. The lasting result is therefore not that regimes always improve PD estimation, but that regime-aware credit models should earn their role through causal, calibration-first, temporally repeated evaluation.
